\documentclass[11pt,letterpaper]{article}
\usepackage[utf8]{inputenc}
\usepackage[T1]{fontenc}
\usepackage[spanish,es-noshorthands,es-tabla]{babel}
\usepackage[margin=2cm,top=2.5cm,bottom=2.5cm]{geometry}
\usepackage{amsmath,amssymb}
\usepackage{enumitem}
\usepackage{fancyhdr}
\usepackage{listings}
\usepackage{xcolor}
\usepackage{tcolorbox}
\usepackage{booktabs}
\usepackage{array}
\raggedbottom
\setlength{\headheight}{14pt}

\pagestyle{fancy}
\fancyhf{}
\fancyhead[L]{Basic C Without Tears}
\fancyhead[C]{Pauta -- Guía 02 Tipo A}
\fancyhead[R]{Stack}
\fancyfoot[L]{Basic-C-Without-Tears}
\fancyfoot[R]{\thepage}
\renewcommand{\headrulewidth}{0.4pt}
\renewcommand{\footrulewidth}{0.4pt}

\lstdefinestyle{customc}{
  language=C,
  numbers=left,
  stepnumber=1,
  numbersep=8pt,
  numberstyle=\tiny\color{gray},
  basicstyle=\footnotesize\ttfamily,
  keywordstyle=\bfseries\color{blue!80!black},
  commentstyle=\itshape\color{green!50!black},
  stringstyle=\color{red!70!black},
  breaklines=true,
  breakatwhitespace=true,
  tabsize=4,
  showstringspaces=false,
  frame=none
}
\lstdefinestyle{pseudocodigo}{
  numbers=none,
  basicstyle=\footnotesize\ttfamily,
  breaklines=true,
  breakatwhitespace=true,
  tabsize=4,
  showstringspaces=false,
  frame=none
}
\newtcolorbox{solucionbox}{
  before=\par\smallskip,
  after=\par\smallskip,
  colback=white,
  colframe=black,
  arc=0mm,
  boxrule=0.6pt,
  left=3mm,
  right=3mm,
  top=2mm,
  bottom=2mm
}

\begin{document}
\begin{center}
    {\LARGE \textbf{Basic C Without Tears}} \\[0.25cm]
    {\Large \textbf{Guía Práctica con Soluciones -- Tipo A -- Guía 02}} \\[0.15cm]
    \small Fundamentos de C, traza de código, strings y matrices \\[0.1cm]
    \textbf{Autor:} Stack
\end{center}
\vspace{0.3cm}

\begin{enumerate}[label=\textbf{\arabic*.}]
    \item \textbf{Fundamentos del Lenguaje (Teoría aplicada).}
    Para los cálculos de tamaño, suponga que un \texttt{char} ocupa 1 byte, un \texttt{int} 4 bytes y un \texttt{double} 8 bytes. Suponga también que \texttt{unsigned char} tiene 8 bits.
    \begin{enumerate}[label=(\alph*)]
        \item Calcule el tamaño de cada objeto declarado y la suma de esos tamaños, en bytes. No considere posibles espacios entre variables.
        \begin{lstlisting}[style=customc,numbers=none]
char etiqueta[6];
int lecturas[3];
double promedio;
        \end{lstlisting}
        \begin{solucionbox}
        \textbf{Solución y Explicación:}\par
        \texttt{etiqueta} ocupa $6\times1=6$ bytes; \texttt{lecturas} ocupa $3\times4=12$ bytes; \texttt{promedio} ocupa 8 bytes. En total son $6+12+8=26$ bytes.
        \end{solucionbox}
        \item Determine si la siguiente expresión es verdadera o falsa en C tradicional e indique su valor entero. Justifique cómo se interpreta el operando negativo.
        \begin{lstlisting}[style=customc,numbers=none]
(0 || -2) && (5 % 2)
        \end{lstlisting}
        \begin{solucionbox}
        \textbf{Solución y Explicación:}\par
        \texttt{0 || -2} vale 1 porque \texttt{-2} es distinto de cero. \texttt{5 \% 2} vale 1. Por tanto, \texttt{1 \&\& 1} es verdadero y su resultado entero es 1.
        En una condición, cero es falso y cualquier valor distinto de cero es verdadero; los operadores lógicos producen exactamente 0 o 1.
        \end{solucionbox}
        \item Si \texttt{dato} es un \texttt{unsigned char} de 8 bits, \textquestiondown{}qué valor queda almacenado después de la asignación? Explique el resultado.
        \begin{lstlisting}[style=customc]
unsigned char dato = 255;
dato = dato + 1;
        \end{lstlisting}
        \begin{solucionbox}
        \textbf{Solución y Explicación:}\par
        El valor almacenado es \textbf{0}. Para la suma, \texttt{dato} se promueve a \texttt{int}, por lo que se calcula $255+1=256$. Al asignarlo a \texttt{unsigned char}, se reduce módulo $256$: el resto es 0. Esta conversión sin signo está definida; no ocurre desbordamiento de \texttt{int}.
        \end{solucionbox}
    \end{enumerate}

    \newpage
    \item \textbf{Análisis y Traza de Código (Matemáticas y Control de Flujo).}
    \textquestiondown{}Qué imprime exactamente este código en pantalla? Suponga que el fragmento está dentro de \texttt{main} y que se incluyó \texttt{stdio.h}. Sin ejecutarlo ni modificarlo, escriba la salida respetando espacios y saltos de línea. Muestre los valores de \texttt{i} y \texttt{total} al terminar cada iteración.
    \begin{lstlisting}[style=customc]
int i = 1;
int total = 0;
while (i <= 4) {
    total += i++ % 3;
    if (i % 2 == 0) {
        printf("P%d ", total);
    } else {
        printf("I%d ", i);
    }
}
switch (total % 3) {
    case 1: printf("uno ");
    case 2: printf("dos ");
    default: printf("fin\n");
}
    \end{lstlisting}
    \begin{solucionbox}
    \textbf{Solución y Explicación:}\par
    En cada iteración, \texttt{i++} aporta el valor anterior de \texttt{i} a la operación y luego incrementa \texttt{i}. Los estados y textos son:
    \begin{center}
    \begin{tabular}{cccc}
    \toprule
    \texttt{i} inicial & \texttt{total} final & \texttt{i} final & Impresión \\
    \midrule
    1 & 1 & 2 & \texttt{P1 } \\
    2 & 3 & 3 & \texttt{I3 } \\
    3 & 3 & 4 & \texttt{P3 } \\
    4 & 4 & 5 & \texttt{I5 } \\
    \bottomrule
    \end{tabular}
    \end{center}
    Al terminar el ciclo, \texttt{total \% 3} vale 1. El \texttt{switch} entra en \texttt{case 1}; como no hay \texttt{break}, continúa por \texttt{case 2} y luego por \texttt{default}.

    \textbf{Salida exacta:}
    \begin{lstlisting}[style=customc,numbers=none]
P1 I3 P3 I5 uno dos fin
    \end{lstlisting}
    Hay un espacio entre cada texto y un único salto de línea después de \texttt{fin}.
    \end{solucionbox}

    \newpage
    \item \textbf{Algoritmia con Strings (Pseudocódigo).}
    Escriba el pseudocódigo de una función que elimine de una cadena \texttt{texto} todas las apariciones del carácter \texttt{objetivo}, modificando la misma cadena. Suponga que \texttt{objetivo} es distinto de \texttt{'\textbackslash 0'}. Recorra los caracteres hasta encontrar ese terminador, conserve el orden de los caracteres restantes y coloque correctamente el terminador al finalizar. No utilice funciones de biblioteca ni un arreglo auxiliar. Describa qué ocurre si la cadena está vacía, si el carácter no aparece o si todos los caracteres coinciden con él.
    \begin{center}
    \begin{tabular}{lll}
    \toprule
    \textbf{Texto} & \textbf{Objetivo} & \textbf{Resultado} \\
    \midrule
    \texttt{"banana"} & \texttt{'a'} & \texttt{"bnn"} \\
    \texttt{"sol"} & \texttt{'x'} & \texttt{"sol"} \\
    \texttt{"aaa"} & \texttt{'a'} & \texttt{""} \\
    \texttt{""} & \texttt{'a'} & \texttt{""} \\
    \bottomrule
    \end{tabular}
    \end{center}
    \begin{solucionbox}
    \textbf{Solución y Explicación:}
    \begin{lstlisting}[style=pseudocodigo]
FUNCION eliminar_caracter(texto, objetivo)
    lectura <- 0
    escritura <- 0

    MIENTRAS texto[lectura] != '\0' HACER
        SI texto[lectura] != objetivo ENTONCES
            texto[escritura] <- texto[lectura]
            escritura <- escritura + 1
        FIN_SI
        lectura <- lectura + 1
    FIN_MIENTRAS

    texto[escritura] <- '\0'
FIN_FUNCION
    \end{lstlisting}
    \texttt{lectura} visita todos los caracteres y \texttt{escritura} solo avanza cuando se conserva uno. Como la escritura nunca adelanta a la lectura, no se sobrescriben caracteres pendientes. El orden se conserva.
    \par Si la cadena está vacía o todos los caracteres se eliminan, se escribe el terminador en la posición 0. Si el objetivo no aparece, la cadena queda igual.
    \end{solucionbox}

    \newpage
    \item \textbf{Algoritmia con Matrices/Arreglos (Pseudocódigo).}
    Una matriz \texttt{notas} tiene \texttt{filas} filas y \texttt{columnas} columnas, ambas cantidades positivas. Escriba pseudocódigo con bucles anidados e índices \texttt{i} y \texttt{j}, comenzando en cero, que calcule e informe el promedio de cada fila. La división debe ser real, sin descartar decimales. Para comprobar la solución, use:
    \[
    \texttt{notas} =
    \begin{bmatrix}
    4 & 8 & 6 & 2 \\
    3 & 5 & 7 & 9 \\
    10 & 0 & 2 & 5
    \end{bmatrix}
    \]
    Indique los tres promedios obtenidos.
    \begin{solucionbox}
    \textbf{Solución y Explicación:}
    \begin{lstlisting}[style=pseudocodigo]
PARA i <- 0 HASTA filas - 1 HACER
    suma <- 0.0
    PARA j <- 0 HASTA columnas - 1 HACER
        suma <- suma + notas[i][j]
    FIN_PARA
    promedio <- suma / columnas   // division real
    ESCRIBIR "Promedio de fila ", i, ": ", promedio
FIN_PARA
    \end{lstlisting}
    Se reinicia \texttt{suma} en cada fila; de lo contrario se acumularían notas de filas anteriores. Los límites incluyen todos los elementos sin salir de la matriz.
    \begin{itemize}[noitemsep]
        \item Fila 0: $(4+8+6+2)/4=5$.
        \item Fila 1: $(3+5+7+9)/4=6$.
        \item Fila 2: $(10+0+2+5)/4=4{,}25$.
    \end{itemize}
    \end{solucionbox}
\end{enumerate}
\end{document}
